\chapter{\Stack}
\label{sec:stack}
\index{\Stack}

\RU{Стек в информатике~--- это одна из наиболее фундаментальных структур данных}%
\EN{The stack is one of the most fundamental data structures in computer science}%
\footnote{\href{http://go.yurichev.com/17119}{wikipedia.org/wiki/Call\_stack}}.\ES{La pila es una de las estructuras de datos más fundamentales en las ciencias de la computación.}

\RU{Технически это просто блок памяти в памяти процесса + регистр \ESP в x86 или \RSP в x64, либо \ac{SP} в ARM, который указывает где-то в пределах этого блока.}
\EN{Technically, it is just a block of memory in process memory along with the \ESP or \RSP register in x86 or x64, or the \ac{SP} register in ARM, as a pointer within that block.}\ES{Técnicamente, es solo un bloque de memoria en la memoria del proceso junto con el registro \ESP o \RSP en x86 o x64, o el registro \ac{SP} en ARM, como un puntero dentro de ese bloque.}

\index{ARM!\Instructions!PUSH}
\index{ARM!\Instructions!POP}
\index{x86!\Instructions!PUSH}
\index{x86!\Instructions!POP}
\RU{Часто используемые инструкции для работы со стеком~--- это \PUSH и \POP (в x86 и Thumb-режиме ARM). 
\PUSH уменьшает \ESP/\RSP/\ac{SP} на 4 в 32-битном режиме (или на 8 в 64-битном),
затем записывает по адресу, на который указывает \ESP/\RSP/\ac{SP}, содержимое своего единственного операнда.}
\EN{The most frequently used stack access instructions are \PUSH and \POP (in both x86 and ARM Thumb-mode). 
\PUSH subtracts from \ESP/\RSP/\ac{SP} 4 in 32-bit mode (or 8 in 64-bit mode) and then writes the contents of its sole operand to the memory address pointed by \ESP/\RSP/\ac{SP}.}\ES{Las instrucciones más frecuentemente utilizadas para acceder a la pila son \PUSH y \POP (tanto en x86 como en el modo Thumb de ARM). \PUSH resta de \ESP/\RSP/\ac{SP} 4 en modo de 32 bits (u 8 en modo de 64 bits) y luego escribe el contenido de su único operando en la dirección de memoria apuntada por \ESP/\RSP/\ac{SP}.}

\RU{\POP это обратная операция~--- сначала достает из \glslink{stack pointer}{указателя стека} значение и помещает его в операнд 
(который очень часто является регистром) и затем увеличивает указатель стека на 4 (или 8).}
\EN{\POP is the reverse operation: retrieve the data from memory pointed by \ac{SP}, 
loads it into the instruction operand (often a register) and then add 4 (or 8) to the \gls{stack pointer}.}\ES{\POP es la operación inversa: recupera los datos de la memoria apuntada por \ac{SP}, los carga en el operando de la instrucción (a menudo un registro) y luego suma 4 (u 8) al \gls{stack pointer}.}

\RU{В самом начале \glslink{stack pointer}{регистр-указатель} указывает на конец стека.}
\EN{After stack allocation, the \gls{stack pointer} points at the bottom of the stack.}\ES{Después de la asignación de la pila, el \gls{stack pointer} apunta al final (fondo) de la pila.}
\RU{\PUSH уменьшает \glslink{stack pointer}{регистр-указатель}, а \POP~--- увеличивает.}
\EN{\PUSH decreases the \gls{stack pointer} and \POP increases it.}\ES{\PUSH disminuye el \gls{stack pointer} y \POP lo incrementa.}
\RU{Конец стека находится в начале блока памяти, выделенного под стек. Это странно, но это так.}
\EN{The bottom of the stack is actually at the beginning of the memory allocated for the stack block. 
It seems strange, but that's the way it is.}\ES{El fondo de la pila está en realidad al principio del bloque de memoria asignado para ella. Parece extraño, pero así es como funciona.}

\ifdefined\IncludeARM
\RU{В процессоре ARM, тем не менее, есть поддержка стеков, растущих как в сторону уменьшения, так и в
сторону увеличения.}
\EN{ARM supports both descending and ascending stacks.}\ES{ARM soporta pilas tanto descendentes como ascendentes.}
\index{ARM!\Instructions!STMFD}
\index{ARM!\Instructions!LDMFD}
\index{ARM!\Instructions!STMED}
\index{ARM!\Instructions!LDMED}
\index{ARM!\Instructions!STMFA}
\index{ARM!\Instructions!LDMFA}
\index{ARM!\Instructions!STMEA}
\index{ARM!\Instructions!LDMEA}

\RU{Например, инструкции}\EN{For example the} 
\ac{STMFD}/\ac{LDMFD}, \ac{STMED}/\ac{LDMED} 
\RU{предназначены для descending-стека 
(растет назад, начиная с высоких адресов в сторону низких).}
\EN{instructions are intended to deal with a descending stack 
(grows downwards, starting with a high address and progressing to a lower one).}\ES{estas instrucciones están destinadas a tratar con una pila descendente (crece hacia abajo, comenzando con una dirección alta y progresando hacia una más baja).}
\RU{Инструкции}\EN{The}
\ac{STMFA}/\ac{LDMFA}, \ac{STMEA}/\ac{LDMEA} 
\RU{предназначены для ascending-стека 
(растет вперед, начиная с низких адресов в сторону высоких).}
\EN{instructions are intended to deal with an ascending stack 
(grows upwards, starting from a low address and progressing to a higher one).}\ES{estas instrucciones están destinadas a tratar con una pila ascendente (crece hacia arriba, comenzando desde una dirección baja y progresando hacia una más alta).}
\fi

\section{\RU{Почему стек растет в обратную сторону?}\EN{Why does the stack grow backwards?}\ES{¿Por qué la pila crece hacia atrás?}}

\RU{Интуитивно мы можем подумать, что, как и любая другая структура данных, стек мог бы расти вперед, 
т.е. в сторону увеличения адресов}
\EN{Intuitively, we might think that the stack grow upwards, i.e. towards
higher addresses, like any other data structure}.\ES{Intuitivamente, podríamos pensar que la pila crece hacia arriba, es decir, hacia direcciones más altas, como cualquier otra estructura de datos.}

\RU{Причина, почему стек растет назад, вероятно, историческая}%
\EN{The reason that the stack grows backward is probably historical}.\ES{La razón por la que la pila crece hacia atrás es probablemente histórica.}
\RU{Когда компьютеры были большие и занимали целую комнату, было очень легко разделить сегмент на две части:
для \glslink{heap}{кучи} и для стека}
\EN{When the computers were big and occupied a whole room, 
it was easy to divide memory into two parts, one for the \gls{heap} and one for the stack}.\ES{Cuando las computadoras eran grandes y ocupaban una habitación entera, era fácil dividir la memoria en dos partes, una para la \gls{heap} (pila de memoria dinámica / montón) y otra para la pila.}
\RU{Заранее было неизвестно, насколько большой может быть \glslink{heap}{куча} или стек, 
так что это решение было самым простым}
\EN{Of course, 
it was unknown how big the \gls{heap} and the stack would be during program execution, 
so this solution was the simplest possible}.\ES{Por supuesto, no se sabía qué tan grandes serían la \gls{heap} y la pila durante la ejecución del programa, por lo que esta solución fue la más simple posible.}

\begin{center}
	\begin{tikzpicture}
	\tikzstyle{every path}=[thick]

	\node [rectangle,draw,minimum width=6cm, minimum height=2cm] (memory) {};
	\node [] [right=0.2cm of memory.west] (heap) {Heap};
	\node [] [left=0.2cm of memory.east] (stack) {Stack};

	\node [] (center1) [right=2cm of memory.west] {};
	\node [] (center2) [left=2cm of memory.east] {};

	\draw [->] (heap) -- (center1);
	\draw [->] (stack) -- (center2);

	\node [] [above left=1.1cm and 0.2cm of heap] (t1) {\RU{Начало кучи}\EN{Start of heap}\ES{Inicio de la heap}};
	\node [] [above right=1.1cm and 0.2cm of stack] (t2) {\RU{Вершина стека}\EN{Start of stack}\ES{Inicio de la pila}};

	\draw [->] (t1) -- (memory.west);
	\draw [->] (t2) -- (memory.east);

	\end{tikzpicture}
\end{center}

\RU{В}\EN{In}\ES{En} \cite{Ritchie74} \RU{можно прочитать}\EN{we can read}\ES{podemos leer}:

\begin{framed}
\begin{quotation}
The user-core part of an image is divided into three logical segments. The program text segment begins at location 0 in the virtual address space. During execution, this segment is write-protected and a single copy of it is shared among all processes executing the same program. At the first 8K byte boundary above the program text segment in the virtual address space begins a nonshared, writable data segment, the size of which may be extended by a system call. Starting at the highest address in the virtual address space is a stack segment, which automatically grows downward as the hardware's stack pointer fluctuates.
\end{quotation}
\end{framed}

\RU{Это немного напоминает как некоторые студенты
пишут два конспеккта в одной тетрадке:
первый конспект начинается обычным образом, второй пишется с конца, перевернув тетрадку.
Конспекты могут встретиться где-то посредине, в случае недостатка свободного места.}
\EN{This reminds us how some students write two lecture notes using only one notebook:
notes for the first lecture are written as usual, 
and notes for the second one are written from the end of notebook, by flipping it.
Notes may meet each other somewhere in between, in case of lack of free space.}\ES{Esto nos recuerda cómo algunos estudiantes escriben dos apuntes de clase usando un solo cuaderno: los apuntes para la primera clase se escriben de manera habitual, y los de la segunda se escriben desde el final del cuaderno, volteándolo. Los apuntes pueden encontrarse en algún lugar intermedio, en caso de falta de espacio libre.}

\section{\RU{Для чего используется стек?}\EN{What is the stack used for?}\ES{¿Para qué se utiliza la pila?}}

% subsections
\subsection{\RU{Сохранение адреса возврата управления}
\EN{Save the function's return address}\ES{Guardar la dirección de retorno de la función}}

\subsubsection{x86}

\index{x86!\Instructions!CALL}
\RU{При вызове другой функции через \CALL сначала в стек записывается адрес, указывающий на место после 
инструкции \CALL, затем делается безусловный переход (почти как \TT{JMP}) на адрес, указанный в операнде.} 
\EN{When calling another function with a \CALL instruction the address of the point exactly after the \CALL instruction is saved 
to the stack and then an unconditional jump to the address in the CALL operand is executed.} 
\ES{Al llamar a otra función con una instrucción \CALL, la dirección del punto exactamente posterior a la instrucción \CALL se guarda en la pila, y luego se realiza un salto incondicional a la dirección en el operando de \CALL.}

\index{x86!\Instructions!PUSH}
\index{x86!\Instructions!JMP}
\RU{\CALL~--- это аналог пары инструкций \TT{PUSH address\_after\_call / JMP}.}
\EN{The \CALL instruction is equivalent to a \TT{PUSH address\_after\_call / JMP operand} instruction pair.}
\ES{La instrucción \CALL es equivalente a un par de instrucciones \TT{PUSH address\_after\_call / JMP operand}.}

\index{x86!\Instructions!RET}
\index{x86!\Instructions!POP}
\RU{\RET вытаскивает из стека значение и передает управление по этому адресу~--- 
это аналог пары инструкций \TT{POP tmp / JMP tmp}.}
\EN{\RET fetches a value from the stack and jumps to it~---that is equivalent to a \TT{POP tmp / JMP tmp} instruction pair.}
\ES{\RET extrae un valor de la pila y salta a él; esto es equivalente a un par de instrucciones \TT{POP tmp / JMP tmp}.}

\index{\Stack!\RU{Переполнение стека}\EN{Stack overflow}\ES{Desbordamiento de pila}}
\index{\Recursion}
\RU{Крайне легко устроить переполнение стека, запустив бесконечную рекурсию:}
\EN{Overflowing the stack is straightforward. Just run eternal recursion:}
\ES{Desbordar la pila es sencillo. Solo ejecute una recursión infinita:}

\begin{lstlisting}
void f()
{
	f();
};
\end{lstlisting}

\RU{MSVC 2008 предупреждает о проблеме:}\EN{MSVC 2008 reports the problem:}\ES{MSVC 2008 reporta el problema:}

\begin{lstlisting}
c:\tmp6>cl ss.cpp /Fass.asm
Microsoft (R) 32-bit C/C++ Optimizing Compiler Version 15.00.21022.08 for 80x86
Copyright (C) Microsoft Corporation.  All rights reserved.

ss.cpp
c:\tmp6\ss.cpp(4) : warning C4717: 'f' : recursive on all control paths, function will cause runtime stack overflow
\end{lstlisting}

\dots \RU{но, тем не менее, создает нужный код}\EN{but generates the right code anyway}\ES{pero genera el código correcto de todos modos}:

\begin{lstlisting}
?f@@YAXXZ PROC						; f
; File c:\tmp6\ss.cpp
; Line 2
	push	ebp
	mov	ebp, esp
; Line 3
	call	?f@@YAXXZ				; f
; Line 4
	pop	ebp
	ret	0
?f@@YAXXZ ENDP						; f
\end{lstlisting}

\dots \RU{причем, если включить оптимизацию (\Ox), то будет даже интереснее, без переполнения стека, 
но работать будет \IT{корректно}\footnote{здесь ирония}:}
\EN{Also if we turn on the compiler optimization (\Ox option) the optimized code will not overflow the stack 
and will work \IT{correctly}\footnote{irony here} instead:}
\ES{Además, si activamos la optimización del compilador (opción \Ox), el código optimizado no desbordará la pila y en su lugar funcionará de manera \IT{correcta}\footnote{aquí hay ironía}:}

\begin{lstlisting}
?f@@YAXXZ PROC						; f
; File c:\tmp6\ss.cpp
; Line 2
$LL3@f:
; Line 3
	jmp	SHORT $LL3@f
?f@@YAXXZ ENDP						; f
\end{lstlisting}

\ifdefined\IncludeGCC
\RU{GCC 4.4.1 генерирует точно такой же код в обоих случаях, хотя и не предупреждает о проблеме.}
\EN{GCC 4.4.1 generates similar code in both cases without, however,  issuing any warning about the problem.}
\ES{GCC 4.4.1 genera un código similar en ambos casos sin emitir, sin embargo, ninguna advertencia sobre el problema.}
\fi

\ifdefined\IncludeARM
\subsubsection{ARM}

\index{ARM!\Registers!Link Register}
\RU{Программы для ARM также используют стек для сохранения \ac{RA}, куда нужно вернуться, но несколько иначе}\EN{ARM
programs also use the stack for saving return addresses, but differently}\ES{Los programas de ARM también utilizan la pila para guardar las direcciones de retorno, pero de manera diferente}.
\RU{Как уже упоминалось в секции}\EN{As mentioned in}\ES{Como se mencionó en} \q{\HelloWorldSectionName}~(\myref{sec:hw_ARM}),
\RU{\ac{RA} записывается в регистр}\EN{the \ac{RA} is saved to the}\ES{la \ac{RA} se guarda en el} \ac{LR} (\gls{link register}).
\RU{Но если есть необходимость вызывать какую-то другую функцию и использовать регистр \ac{LR} ещё
раз, его значение желательно сохранить}%
\EN{If one needs, however, to call another function and use the \ac{LR} register
one more time its value has to be saved}\ES{Sin embargo, si se necesita llamar a otra función y usar el registro \ac{LR} una vez más, su valor debe ser guardado}.
\index{Function prologue}
\RU{Обычно это происходит в прологе функции, часто мы видим там инструкцию вроде}
\EN{Usually it is saved in the function prologue. Often, we see instructions like}\ES{Por lo general, esto ocurre en el prólogo de la función. A menudo, vemos instrucciones como}
\index{ARM!\Instructions!PUSH}
\index{ARM!\Instructions!POP}
\TT{PUSH {R4-R7,LR}} \RU{, а в эпилоге}\EN{along with this instruction in epilogue}\ES{, y en el epílogo}
\TT{POP {R4-R7,PC}}\RU{~--- так сохраняются регистры, которые будут использоваться в текущей функции, в том числе}
\EN{---thus register values
to be used in the function are saved in the stack, including}\ES{; de este modo se guardan en la pila los valores de los registros que se utilizarán en la función, incluyendo} \ac{LR}.

\index{ARM!Leaf function}
\RU{Тем не менее, если некая функция не вызывает никаких более функций, в терминологии \ac{RISC} она называется}
\EN{Nevertheless, if a function never calls any other function, in \ac{RISC} terminology it is called a}\ES{Sin embargo, si una función nunca llama a ninguna otra función, en la terminología \ac{RISC} se le llama una}
\IT{\gls{leaf function}}\footnote{\href{http://go.yurichev.com/17064}{infocenter.arm.com/help/index.jsp?topic=/com.arm.doc.faqs/ka13785.html}}. 
\RU{Как следствие, \q{leaf}-функция не сохраняет регистр \ac{LR} (потому что не изменяет его).}
\EN{As a consequence, leaf functions do not save the \ac{LR} register (because they don't modify it).}\ES{Como consecuencia, las funciones hoja (\IT{leaf functions}) no guardan el registro \ac{LR} (porque no lo modifican).}
\RU{А если эта функция небольшая, использует мало регистров, она может не использовать стек вообще}%
\EN{If such function is small and uses a small number of registers, it may not use the stack at all}\ES{Si dicha función es pequeña y utiliza pocos registros, es posible que no utilice la pila en absoluto}.
\RU{Таким образом, в ARM возможен вызов небольших leaf-функций не используя стек.}
\EN{Thus, it is possible to call leaf functions without using the stack}\ES{Por lo tanto, es posible llamar a funciones hoja sin utilizar la pila}
\RU{Это может быть быстрее чем в старых x86, ведь внешняя память для стека не используется}%
\EN{which can be faster than on older x86 because external RAM is not used for the stack}\ES{, lo cual puede ser más rápido que en los antiguos x86 debido a que no se utiliza la memoria RAM externa para la pila}%
\footnote{\RU{Когда-то, очень давно, на PDP-11 и VAX на инструкцию CALL (вызов других функций) могло тратиться
вплоть до 50\% времени (возможно из-за работы с памятью),
поэтому считалось, что много небольших функций это \glslink{anti-pattern}{анти-паттерн}}%
\EN{Some time ago, on PDP-11 and VAX, the CALL instruction (calling other functions) was expensive; up to 50\%
of execution time might be spent on it, so it was considered that having a big number of small functions is an \gls{anti-pattern}}\ES{Hace algún tiempo, en PDP-11 y VAX, la instrucción CALL (llamar a otras funciones) era costosa; se podía gastar hasta el 50\% del tiempo de ejecución en ella, por lo que se consideraba que tener un gran número de funciones pequeñas era un \gls{anti-pattern}} \cite[Chapter 4, Part II]{Raymond:2003:AUP:829549}.}.
\RU{Либо это может быть полезным для тех ситуаций, когда память для стека ещё не выделена, либо недоступна}%
\EN{This can be also useful for situations when memory for the stack is not yet allocated or not available}\ES{Esto también puede ser útil para situaciones en las que la memoria para la pila aún no ha sido asignada o no está disponible}.

\EN{Some examples of leaf functions:}\RU{Некоторые примеры таких функций:}\ES{Algunos ejemplos de funciones hoja:}
\myref{ARM_leaf_example1}, \myref{ARM_leaf_example2}, 
\myref{ARM_leaf_example3}, \myref{ARM_leaf_example4}, \myref{ARM_leaf_example5},
\myref{ARM_leaf_example6}, \myref{ARM_leaf_example7}, \myref{ARM_leaf_example10}.
\fi

\subsection{\RU{Передача параметров функции}\EN{Passing function arguments}\ES{Paso de argumentos de la función}}

\RU{Самый распространенный способ передачи параметров в x86 называется}
\EN{The most popular way to pass parameters in x86 is called} \q{cdecl}:
\ES{La forma más común de pasar parámetros en x86 se llama} \q{cdecl}:

\begin{lstlisting}
push arg3
push arg2
push arg1
call f
add esp, 12 ; 4*3=12
\end{lstlisting}

\RU{Вызываемая функция получает свои параметры также через указатель стека.}
\EN{\Gls{callee} functions get their arguments via the stack pointer.}
\ES{Las funciones llamadas (\gls{callee}) obtienen sus argumentos a través del puntero de pila.}

\RU{Следовательно, так расположены значения в стеке перед исполнением самой первой инструкции
функции \ttf{}:}
\EN{Therefore, this is how the argument values are located in the stack before the execution
of the \ttf{} function's very first instruction:}
\ES{Por lo tanto, así es como se ubican los valores de los argumentos en la pila antes de la ejecución de la primera instrucción de la función \ttf{}:}

\begin{center}
\begin{tabular}{ | l | l | }
\hline
ESP & \RU{адрес возврата}\EN{return address}\ES{dirección de retorno} \\
\hline
ESP+4 & \argument \#1, \MarkedInIDAAs{} \TT{arg\_0} \\
\hline
ESP+8 & \argument \#2, \MarkedInIDAAs{} \TT{arg\_4} \\
\hline
ESP+0xC & \argument \#3, \MarkedInIDAAs{} \TT{arg\_8} \\
\hline
\dots & \dots \\
\hline
\end{tabular}
\end{center}

\ifx\LITE\undefined
\RU{См. также в соответствующем разделе о других способах передачи аргументов через стек}
\EN{For more information on other calling conventions see also section}~(\myref{sec:callingconventions}).
\ES{Para obtener más información sobre otras convenciones de llamada, consulte también la sección}~(\myref{sec:callingconventions}).
\fi
\RU{Важно отметить, что, в общем, никто не заставляет программистов передавать параметры именно через стек,
это не является требованием к исполняемому коду.}
\EN{It is worth noting that nothing obliges programmers to pass arguments through the stack. It is not a requirement.}
\ES{Vale la pena señalar que nada obliga a los programadores a pasar argumentos a través de la pila. No es un requisito.}
\RU{Вы можете делать это совершенно иначе, не используя стек вообще.}
\EN{One could implement any other method without using the stack at all.}
\ES{Se podría implementar cualquier otro método sin usar la pila en absoluto.}

\RU{К примеру, можно выделять в \glslink{heap}{куче} место для аргументов, 
заполнять их и передавать в функцию указатель на это место через \EAX. И это вполне будет работать}%
\EN{For example, it is possible to allocate a space for arguments in the \gls{heap}, fill it and pass it to a function 
via a pointer to this block in the \EAX register. This will work}%
\ES{Por ejemplo, es posible asignar espacio para los argumentos en el montículo (\gls{heap}), llenarlo y pasarlo a una función a través de un puntero a este bloque en el registro \EAX. Esto funcionará}%
\footnote{\RU{Например, в книге Дональда Кнута \q{Искусство программирования}, в разделе 1.4.1 
посвященном подпрограммам \cite[раздел 1.4.1]{Knuth:1998:ACP:521463}, 
мы можем прочитать о возможности располагать параметры для вызываемой подпрограммы после инструкции \JMP,
передающей управление подпрограмме. Кнут описывает, что это было особенно удобно для компьютеров IBM System/360.}%
\EN{For example, in the \q{The Art of Computer Programming} book by Donald Knuth, 
in section 1.4.1 dedicated to subroutines \cite[section 1.4.1]{Knuth:1998:ACP:521463},
we could read that one way to supply arguments to a subroutine is simply to list them after the \JMP instruction
passing control to subroutine. Knuth explains that this method was particularly convenient on IBM System/360.}
\ES{Por ejemplo, en el libro \q{The Art of Computer Programming} de Donald Knuth, en la sección 1.4.1 dedicada a subrutinas \cite[sección 1.4.1]{Knuth:1998:ACP:521463}, podemos leer que una forma de proporcionar argumentos a una subrutina es simplemente enumerarlos después de la instrucción \JMP que pasa el control a la subrutina. Knuth explica que este método era particularmente conveniente en la IBM System/360.}}.
\RU{Однако традиционно сложилось, что в x86 и ARM передача аргументов происходит именно через стек.}
\EN{However, it is a convenient custom in x86 and ARM to use the stack for this purpose.} \\
\ES{Sin embargo, es una costumbre conveniente en x86 y ARM utilizar la pila para este propósito.} \\

\RU{Кстати, вызываемая функция не имеет информации о количестве переданных ей аргументов.}
\EN{By the way, the \gls{callee} function does not have any information about how many arguments were passed.}
\ES{Por cierto, la función llamada (\gls{callee}) no tiene ninguna información sobre cuántos argumentos le fueron pasados.}
\RU{Функции Си с переменным количеством аргументов (как \printf) определяют их количество по 
спецификаторам строки формата (начинающиеся со знака \%).}
\EN{C functions with a variable number of arguments (like \printf) determine their number using format string  specifiers (which begin with the \% symbol).}
\ES{Las funciones de C con un número variable de argumentos (como \printf) determinan su cantidad utilizando los especificadores de la cadena de formato (que comienzan con el símbolo \%).}
\RU{Если написать что-то вроде}\EN{If we write something like}\ES{Si escribimos algo como} 

\begin{lstlisting}
printf("%d %d %d", 1234);
\end{lstlisting}

\printf \RU{выведет 1234, затем ещё два случайных числа, которые волею случая оказались в стеке рядом.}
\EN{will print 1234, and then two random numbers, which were laying next to it in the stack.}
\ES{imprimirá 1234 y luego dos números aleatorios que se encontraban junto a él en la pila.}\\

\RU{Вот почему не так уж и важно, как объявлять функцию \main}
\EN{That's why it is not very important how we declare the \main function}: \RU{как}\EN{as}\ES{como} \main, 
\TT{main(int argc, char *argv[])} 
\RU{либо}\EN{or}\ES{o} \TT{main(int argc, char *argv[], char *envp[])}.

\RU{В реальности, \ac{CRT}-код вызывает \main примерно так:}
\EN{In fact, the \ac{CRT}-code is calling \main roughly as:}
\ES{De hecho, el código de \ac{CRT} llama a \main aproximadamente como:}

\begin{lstlisting}
push envp
push argv
push argc
call main
...
\end{lstlisting}

\RU{Если вы объявляете \main без аргументов, они, тем не менее, присутствуют в стеке, но не используются.}
\EN{If you declare \main as \main without arguments, they are, nevertheless, still present in the stack, but
are not used.}
\ES{Si declara \main sin argumentos, estos, no obstante, siguen estando presentes en la pila, pero no se utilizan.}
\RU{Если вы объявите \main как}\EN{If you declare \main as}\ES{Si declara \main como} \TT{main(int argc, char *argv[])}, 
\RU{вы можете использовать два первых аргумента, а третий останется для вашей функции \q{невидимым}.}
\EN{you will be able to use first two arguments, and the third will remain \q{invisible} for your function.}
\ES{podrá usar los dos primeros argumentos y el tercero permanecerá \q{invisible} para su función.}
\RU{Более того, можно даже объявить}\EN{Even more, it is possible to declare}\ES{Es más, incluso es posible declarar} \TT{main(int argc)}, 
\RU{и это будет работать}\EN{and it will work}\ES{y funcionará}.


\subsection{\RU{Хранение локальных переменных}\EN{Local variable storage}\ES{Almacenamiento de variables locales}}

\RU{Функция может выделить для себя некоторое место в стеке для локальных переменных, просто отодвинув 
\glslink{stack pointer}{указатель стека} глубже к концу стека.}
\EN{A function could allocate space in the stack for its local variables just by shifting 
the \gls{stack pointer} towards the stack bottom.}
\ES{Una función puede asignar espacio en la pila para sus variables locales simplemente desplazando el \glslink{stack pointer}{puntero de pila} hacia el final de la pila.}
\RU{Это очень быстро вне зависимости от количества локальных переменных.}
\EN{Hence, it's very fast, no matter how many local variables are defined.}
\ES{Por lo tanto, es muy rápido, sin importar cuántas variables locales estén definidas.}

\RU{Хранить локальные переменные в стеке не является необходимым требованием. 
Вы можете хранить локальные переменные где угодно. 
Но по традиции всё сложилось так.}
\EN{It is also not a requirement to store local variables in the stack.
You could store local variables wherever you like, 
but traditionally this is how it's done.}
\ES{Tampoco es un requisito almacenar las variables locales en la pila. Podría almacenar variables locales donde quiera, pero tradicionalmente es así como se hace.}

\subsection{x86: \RU{Функция alloca()}\EN{alloca() function}\ES{Función alloca()}}
\label{alloca}
\index{\CStandardLibrary!alloca()}
\RU{Интересен случай с функцией \TT{alloca()}}%
\EN{It is worth noting the \TT{alloca()} function}\ES{Vale la pena notar la función \TT{alloca()}}\footnote{
\RU{В MSVC, реализацию функции можно посмотреть в файлах}%
\EN{In MSVC, the function implementation can be found in} 
\ES{En MSVC, la implementación de la función se puede encontrar en} 
  \TT{alloca16.asm} 
  \AndENRU 
  \TT{chkstk.asm} 
  \InENRU 
  \TT{C:\textbackslash{}Program Files (x86)\textbackslash{}Microsoft Visual Studio 10.0\textbackslash{}VC\textbackslash{}crt\textbackslash{}src\textbackslash{}intel}}. 

\RU{Эта функция работает как \TT{malloc()}, но выделяет память прямо в стеке.} 
\EN{This function works like \TT{malloc()}, but allocates memory directly on the stack.}
\ES{Esta función funciona como \TT{malloc()}, pero asigna memoria directamente en la pila.}

\RU{Память освобождать через \TT{free()} не нужно, так как эпилог функции~(\myref{sec:prologepilog})
вернет \ESP в изначальное состояние и выделенная память просто \IT{выкидывается}.}
\EN{The allocated memory chunk does not need to be freed via a \TT{free()} function call, since the 
function epilogue~(\myref{sec:prologepilog}) returns \ESP back to its initial state and 
the allocated memory is just \IT{dropped}.}
\ES{El bloque de memoria asignado no necesita ser liberado a través de una llamada a la función \TT{free()}, ya que el epílogo de la función~(\myref{sec:prologepilog}) devuelve \ESP a su estado inicial y la memoria asignada simplemente se \IT{desecha}.}

\RU{Интересна реализация функции \TT{alloca()}.}
\EN{It is worth noting how \TT{alloca()} is implemented.}
\ES{Vale la pena señalar cómo está implementada \TT{alloca()}.}

\RU{Эта функция, если упрощенно, просто сдвигает \ESP вглубь стека 
на столько байт, сколько вам нужно и возвращает \ESP в качестве указателя на выделенный блок.}
\EN{In simple terms, this function just shifts \ESP downwards toward the stack bottom by the number of bytes you need and sets \ESP as a pointer to the \IT{allocated} block.}
\ES{En términos sencillos, esta función solo desplaza \ESP hacia abajo, hacia el fondo de la pila, por la cantidad de bytes que necesite, y devuelve \ESP como un puntero al bloque \IT{asignado}.}
\RU{Попробуем:}\EN{Let's try:}\ES{Probemos:}

\lstinputlisting{patterns/02_stack/04_alloca/2_1.c}

\RU{Функция \TT{\_snprintf()} работает так же, как и \printf, только вместо выдачи результата в \gls{stdout} (т.е. на терминал или в консоль),
записывает его в буфер \TT{buf}. Функция \puts выдает содержимое буфера \TT{buf} в \gls{stdout}. Конечно, можно было бы
заменить оба этих вызова на один \printf, но здесь нужно проиллюстрировать использование небольшого буфера.}%
\EN{\TT{\_snprintf()} function works just like \printf, but instead of dumping the result into \gls{stdout} (e.g., to terminal or 
console), it writes it to the \TT{buf} buffer. Function \puts copies \TT{buf} contents to \gls{stdout}. Of course, these two
function calls might be replaced by one \printf call, but we have to illustrate small buffer usage.}
\ES{La función \TT{\_snprintf()} funciona igual que \printf, pero en lugar de enviar el resultado a \gls{stdout} (por ejemplo, al terminal o consola), lo escribe en el búfer \TT{buf}. La función \puts copia el contenido de \TT{buf} a \gls{stdout}. Por supuesto, estas dos llamadas a funciones podrían reemplazarse por una sola llamada a \printf, pero aquí tenemos que ilustrar el uso de un búfer pequeño.}%

\subsubsection{MSVC}

\RU{Компилируем}\EN{Let's compile}\ES{Compilemos} (MSVC 2010):

\lstinputlisting[caption=MSVC 2010]{patterns/02_stack/04_alloca/2_2_msvc.asm}

\index{Compiler intrinsic}
\RU {Единственный параметр в \TT{alloca()} передается через \EAX, а не как обычно через стек}%
\EN{The sole \TT{alloca()} argument is passed via \EAX (instead of pushing it into the stack)}%
\ES{El único argumento de \TT{alloca()} se pasa a través de \EAX (en lugar de empujarlo a la pila)}%
\footnote{
\RU{Это потому, что alloca()~--- это не сколько функция, сколько т.н. \IT{compiler intrinsic}%
\ifx\LITE\undefined
(\myref{sec:compiler_intrinsic})
\fi
}%
\EN{It is because alloca() is rather a compiler intrinsic 
\ifx\LITE\undefined
(\myref{sec:compiler_intrinsic}) 
\fi
than a normal function}.
\ES{Esto se debe a que alloca() es más bien una función intrínseca del compilador (\IT{compiler intrinsic})\ifx\LITE\undefined (\myref{sec:compiler_intrinsic})\fi en lugar de una función normal}.

\RU{Одна из причин, почему здесь нужна именно функция, а не несколько инструкций прямо в коде в том, что в реализации 
функции alloca() от \ac{MSVC}
есть также код, читающий из только что выделенной памяти, чтобы \ac{OS} подключила физическую память к этому региону \ac{VM}.}
\EN{One of the reasons we need a separate function instead of just a couple of instructions in the code,
is because the \ac{MSVC} alloca() implementation also has code which reads from the memory just allocated, in order to let the \ac{OS} map
physical memory to this \ac{VM} region.}
\ES{Una de las razones por las que necesitamos una función separada, en lugar de solo un par de instrucciones en el código, es porque la implementación de alloca() de \ac{MSVC} también contiene código que lee de la memoria recién asignada, con el fin de permitir que el \ac{OS} asigne memoria física a esta región de la \ac{VM}.}
}.
\RU{После вызова \TT{alloca()} \ESP указывает на блок в 600 байт, который 
мы можем использовать под \TT{buf}.}
\EN{After the \TT{alloca()} call, \ESP points to the block of 600 bytes and we can 
use it as memory for the \TT{buf} array.}
\ES{Después de la llamada a \TT{alloca()}, \ESP apunta a un bloque de 600 bytes y podemos usarlo como memoria para el búfer \TT{buf}.}

\ifdefined\IncludeGCC
\subsubsection{GCC + \IntelSyntax}

\RU{А GCC 4.4.1 обходится без вызова других функций:}
\EN{GCC 4.4.1 does the same without calling external functions:}
\ES{Y GCC 4.4.1 hace lo mismo sin llamar a funciones externas:}

\lstinputlisting[caption=GCC 4.7.3]{patterns/02_stack/04_alloca/2_1_gcc_intel_O3.asm.\LANG}

\subsubsection{GCC + \ATTSyntax}

\RU{Посмотрим на тот же код, только в синтаксисе AT\&T}\EN{Let's see the same code, but in AT\&T syntax}\ES{Veamos el mismo código, pero en sintaxis AT\&T}:

\lstinputlisting[caption=GCC 4.7.3]{patterns/02_stack/04_alloca/2_1_gcc_ATT_O3.s}

\index{\ATTSyntax}
\RU{Всё то же самое, что и в прошлом листинге.}\EN{The code is the same as in the previous listing.}\ES{El código es el mismo que en el listado anterior.}

\RU{Кстати}\EN{By the way}\ES{Por cierto}, \TT{movl \$3, 20(\%esp)}%
\RU{~--- это аналог}\EN{corresponds to}\ES{corresponde a} \TT{mov DWORD PTR [esp+20], 3}
\RU{в синтаксисе Intel.}\EN{ in Intel-syntax.}\ES{ en sintaxis Intel.}
\RU{Адресация памяти в виде \IT{регистр+смещение} записывается в синтаксисе AT\&T как \TT{смещение(\%{регистр})}.}
\EN{In AT\&T syntax register+offset format of addressing memory looks like \TT{offset(\%{register})}.}
\ES{En la sintaxis AT\&T, el formato de direccionamiento de memoria de \IT{registro+desplazamiento} se escribe como \TT{desplazamiento(\%{registro})}.}
\fi

\subsection{(Windows) SEH}
\index{Windows!Structured Exception Handling}

\RU{В стеке хранятся записи \ac{SEH} для функции (если они присутствуют)}%
\EN{\ac{SEH} records are also stored on the stack (if they present).}%
\ES{Los registros de \ac{SEH} también se almacenan en la pila (si están presentes)}.

\ifx\LITE\undefined
\RU{Читайте больше о нем здесь}\EN{Read more about it}\ES{Lea más al respecto}: (\myref{sec:SEH}).
\fi

\subsection{\RU{Защита от переполнений буфера}\EN{Buffer overflow protection}\ES{Protección contra desbordamientos de búfer}}

\RU{Здесь больше об этом}\EN{More about it here}\ES{Más sobre esto aquí}~(\myref{subsec:bufferoverflow}).


\subsection{\EN{Automatic deallocation of data in stack}\RU{Автоматическое освобождение данных в стеке}\ES{Liberación automática de datos en la pila}}

\RU{Возможно, причина хранения локальных переменных и SEH-записей в стеке в том, что после выхода из функции, всё эти данные освобождаются автоматически,
используя только одну инструкцию корректирования указателя стека (часто это ADD).}
\EN{Perhaps, the reason of storing local variables and SEH records in stack is because they are freed automatically upon function exit,
using just one instruction of stack pointer correcting (it is often ADD).}\ES{Tal vez, la razón para almacenar variables locales y registros SEH en la pila sea porque se liberan automáticamente al salir de la función, usando solo una instrucción de corrección del puntero de pila (a menudo es ADD).}
\RU{Аргументы функций, можно сказать, тоже освобождаются автоматически в конце функции.}
\EN{Function arguments, as we could say, also deallocated automatically at the end of function.}\ES{Los argumentos de las funciones, por así decirlo, también se liberan automáticamente al final de la función.}
\RU{А всё что хранится в куче (\IT{heap}) нужно освобождать явно.}
\EN{On contrary, everything stored in \IT{heap} must be deallocated explicitly.}\ES{Por el contrario, todo lo almacenado en la \IT{heap} debe ser liberado explíquitamente.}

% sections
\section{\RU{Разметка типичного стека}\EN{Typical stack layout}\ES{Diseño típico de la pila}}

\RU{Разметка типичного стека в 32-битной среде
перед исполнением самой первой инструкции функции выглядит так:}
\EN{A typical stack layout in a 32-bit environment at the start of a function, 
before the first instruction execution looks like this:}
\ES{Un diseño típico de la pila en un entorno de 32 bits al inicio de una función, antes de la ejecución de la primera instrucción, se ve así:}

\begin{center}
\begin{tabular}{ | l | l | }
\hline
\dots & \dots \\
\hline
ESP-0xC & \RU{локальная переменная}\EN{local variable}\ES{variable local} \#2, \MarkedInIDAAs{} \TT{var\_8} \\
\hline
ESP-8 & \RU{локальная переменная}\EN{local variable}\ES{variable local} \#1, \MarkedInIDAAs{} \TT{var\_4} \\
\hline
ESP-4 & \RU{сохраненное значение}\EN{saved value of}\ES{valor guardado de} \EBP \\
\hline
ESP & \RU{адрес возврата}\EN{return address}\ES{dirección de retorno} \\
\hline
ESP+4 & \argument \#1, \MarkedInIDAAs{} \TT{arg\_0} \\
\hline
ESP+8 & \argument \#2, \MarkedInIDAAs{} \TT{arg\_4} \\
\hline
ESP+0xC & \argument \#3, \MarkedInIDAAs{} \TT{arg\_8} \\
\hline
\dots & \dots \\
\hline
\end{tabular}
\end{center}

\ifx\LITE\undefined
\section{\RU{Мусор в стеке}\EN{Noise in stack}\ES{Ruido en la pila}}

\RU{Часто в этой книге говорится о \q{шуме} или \q{мусоре} в стеке или памяти.}
\EN{Often in this book \q{noise} or \q{garbage} values in stack or memory are mentioned.}
\ES{A menudo en este libro se mencionan valores de \q{ruido} o \q{basura} en la pila o en la memoria.}
\RU{Откуда он берется}\EN{Where do they come from}\ES{¿De dónde vienen}?
\RU{Это то, что осталось там после исполнения предыдущих функций.}
\EN{These are what was left in there after other functions' executions.}
\ES{Esto es lo que quedó allí después de las ejecuciones de otras funciones.}
\RU{Короткий пример}\EN{Short example}\ES{Ejemplo corto}:

\lstinputlisting{patterns/02_stack/08_noise/st.c}

\RU{Компилируем}\EN{Compiling}\ES{Compilando}\dots

\lstinputlisting[caption=\NonOptimizing MSVC 2010]{patterns/02_stack/08_noise/st.asm}

\RU{Компилятор поворчит немного}\EN{The compiler will grumble a little bit}\ES{El compilador se quejará un poco}\dots

\begin{lstlisting}
c:\Polygon\c>cl st.c /Fast.asm /MD
Microsoft (R) 32-bit C/C++ Optimizing Compiler Version 16.00.40219.01 for 80x86
Copyright (C) Microsoft Corporation.  All rights reserved.

st.c
c:\polygon\c\st.c(11) : warning C4700: uninitialized local variable 'c' used
c:\polygon\c\st.c(11) : warning C4700: uninitialized local variable 'b' used
c:\polygon\c\st.c(11) : warning C4700: uninitialized local variable 'a' used
Microsoft (R) Incremental Linker Version 10.00.40219.01
Copyright (C) Microsoft Corporation.  All rights reserved.

/out:st.exe
st.obj
\end{lstlisting}

\RU{Но когда мы запускаем}\EN{But when we run the compiled program}\ES{Pero cuando ejecutamos el programa compilado}\dots

\begin{lstlisting}
c:\Polygon\c>st
1, 2, 3
\end{lstlisting}

\RU{Ох. Вот это странно. Мы ведь не устанавливали значения никаких переменных в}\EN{Oh, 
what a weird thing! We did not set any variables in}\ES{¡Oh, qué cosa tan extraña! No establecimos ninguna variable en} \TT{f2()}. 
\RU{Эти значения --- это \q{привидения}, которые всё ещё в стеке.}
\EN{These are \q{ghosts} values, which are still in the stack.}
\ES{Estos son valores \q{fantasma} que todavía están en la pila.}

\clearpage
\RU{Загрузим пример в}\EN{Let's load the example into}\ES{Carguemos el ejemplo en} \olly:

\begin{figure}[H]
\centering
\includegraphics[scale=\FigScale]{patterns/02_stack/08_noise/olly1.png}
\caption{\olly: \TT{f1()}}
\label{fig:stack_noise_olly1}
\end{figure}

\RU{Когда}\EN{When}\ES{Cuando} \TT{f1()} \RU{заполняет переменные}\EN{assigns the variables}\ES{asigna las variables} $a$, $b$ \AndENRU $c$ 
\RU{они сохраняются по адресу}\EN{, their values are stored at the address}\ES{sus valores se almacenan en la dirección} \TT{0x14F860} 
\RU{\etc{}.}\EN{and so on.}\ES{ y así sucesivamente.}

\clearpage
\RU{А когда исполняется}\EN{And when}\ES{Y cuando} \TT{f2()}\EN{ executes}\ES{ se ejecuta}:

\begin{figure}[H]
\centering
\includegraphics[scale=\FigScale]{patterns/02_stack/08_noise/olly2.png}
\caption{\olly: \TT{f2()}}
\label{fig:stack_noise_olly2}
\end{figure}

... $a$, $b$ \AndENRU $c$ \RU{в функции}\EN{of}\ES{de} \TT{f2()} \RU{находятся по тем же адресам!}
\EN{are located at the same addresses!}
\ES{¡se encuentran en las mismas direcciones!}
\RU{Пока никто не перезаписал их, так что они здесь в нетронутом виде.}
\EN{No one has overwritten the values yet, so at that point they are still untouched.}
\ES{Nadie ha sobrescrito los valores todavía, por lo que en ese momento aún están intactos.}

\RU{Для создания такой странной ситуации несколько функций должны исполняться друг за другом
и \ac{SP} должен быть одинаковым при входе в функции, т.е. у функций должно быть равное количество
аргументов). Тогда локальные переменные будут расположены в том же месте стека.}
\EN{So, for this weird situation, several functions has to be called one after another and
\ac{SP} has to be the same at each function entry (i.e., they  has same number
of arguments). Then the local variables will be located at the same positions in the stack.}
\ES{Por lo tanto, para esta situación extraña, se tienen que llamar varias funciones una tras otra y el \ac{SP} tiene que ser el mismo en cada entrada de la función (es decir, deben tener la misma cantidad de argumentos). Entonces las variables locales se ubicarán en las mismas posiciones en la pila.}

\RU{Подводя итоги, все значения в стеке (да и памяти вообще) это значения оставшиеся от 
исполнения предыдущих функций.}
\EN{Summarizing, all values in the stack (and memory cells in general) 
have values left there from previous function executions.}
\ES{Resumiendo, todos los valores en la pila (y las celdas de memoria en general) tienen valores que quedaron allí de ejecuciones de funciones anteriores.}
\RU{Строго говоря, они не случайны, они скорее непредсказуемы.}
\EN{They are not random in the strict sense, but rather have unpredictable values.}
\ES{No son aleatorios en el sentido estricto, sino que tienen valores impredecibles.}

\RU{А как иначе}\EN{Is there other option}\ES{¿Hay otra opción}?
\RU{Можно было бы очищать части стека перед исполнением каждой функции,
но это слишком много лишней (и ненужной) работы.}
\EN{It probably would be possible to clear portions of the stack before each function execution,
but that's too much extra (and unnecessary) work.}
\ES{Probablemente sería posible limpiar porciones de la pila antes de cada ejecución de función, pero eso es demasiado trabajo adicional (e innecesario).}

\fi
\ifdefined\IncludeExercises
\input{patterns/02_stack/exercises}
\fi
